El presente proyecto nace en un contexto muy concreto: la \textbf{distribución de productos profesionales de peluquería}. Aunque desde fuera pueda parecer una actividad puramente comercial, en la práctica implica coordinar de forma continua pedidos, clientes, visitas, cobros, entregas y un catálogo de productos que debe mantenerse actualizado. En mi caso, el contacto con esta problemática surgió a partir de una propuesta que recibí mientras buscaba temática para el \textbf{Trabajo Fin de Grado}. Esta oportunidad me permitió conocer de primera mano la forma de trabajar del distribuidor y comprobar cómo la \textit{falta de trazabilidad} y la \textit{dependencia de procesos manuales} dificultaban su actividad diaria.

Durante el análisis inicial se observa que muchas de estas tareas se apoyan en herramientas separadas entre sí, como \textit{hojas de cálculo}, \textit{bases de datos locales} o \textit{aplicaciones independientes}. El problema no reside únicamente en el uso de varias herramientas, sino en la fragmentación de la información, lo que dificulta consultarla de forma rápida, mantenerla actualizada y seguir con claridad la actividad diaria. Un ejemplo claro de esto era la necesidad de consultar datos de clientes, pedidos o cobros en documentos distintos, teniendo que contrastar manualmente la información antes de continuar con una gestión.

En el caso estudiado, esta forma de trabajo genera una \textbf{carga administrativa elevada} y una dependencia de procesos manuales. Esto limita la agilidad del distribuidor, especialmente cuando necesita consultar información de clientes, revisar pedidos pendientes, organizar visitas o tener una visión general del estado del negocio. Además, la relación digital con las peluquerías profesionales todavía es reducida, lo que impide aprovechar mejor la información que se genera en la actividad comercial y dificulta ofrecer una atención más ordenada y personalizada.

Por tanto, la motivación principal surge de la posibilidad de aplicar los conocimientos obtenidos durante el Grado en Ingeniería Informática a un \textbf{problema real y cercano}, trabajando sobre un caso con \textbf{necesidades reales} en el que cada decisión de diseño tiene una utilidad práctica. El objetivo no es desarrollar una aplicación web de forma aislada, sino plantear una herramienta que pueda integrar en una única plataforma las funciones principales de gestión del distribuidor y que sirva como punto de partida para \textit{digitalizar progresivamente} su forma de trabajar. Con ello, se busca mejorar la organización interna, facilitar el acceso a la información y ofrecer a los clientes profesionales un entorno digital propio que refuerce su relación con el distribuidor. Además, el desarrollo de la aplicación permite aplicar de forma conjunta competencias relacionadas con el \textit{análisis de requisitos}, el \textit{diseño de sistemas de información}, las \textit{bases de datos}, el \textit{desarrollo web}, la \textit{autenticación}, el \textit{control de roles} y el \textit{despliegue de aplicaciones}.